\section{Bandidos Contextuales Lineales y su Relación con RLHF}

\subsection*{a) Recompensas estimadas y bonos de incertidumbre}

LinUCB usa la predicción lineal $x_t^\top\hat\theta_a$ y un bono que
depende de la incertidumbre en la dirección de
$x_t$~\cite{li2010}. Con $x_t=(1,2)^\top$, las predicciones son
\begin{align*}
    \hat r_{t,1}&=1(0.5)+2(0.4)=1.3,\\
    \hat r_{t,2}&=1(0.4)+2(0.3)=1.0.
\end{align*}
Las matrices son diagonales y $\alpha=1$, por lo que los bonos son
\begin{align*}
    u_{t,1}&=\sqrt{0.2(1)^2+0.1(2)^2}
      =\sqrt{0.6}\approx 0.775,\\
    u_{t,2}&=\sqrt{0.5(1)^2+0.5(2)^2}
      =\sqrt{2.5}\approx 1.581.
\end{align*}

\subsection*{b) Selección y actualización tras el clic}

Los índices son $1.3+0.775=2.075$ para $a_1$ y
$1.0+1.581=2.581$ para $a_2$. LinUCB selecciona $a_2$ porque su
bono de incertidumbre compensa su menor predicción de recompensa.
El clic aporta $r_t=1$ y actualiza las estadísticas de regresión
del anuncio seleccionado,
\begin{align*}
    A_2^{\mathrm{nuevo}}
    &=A_2+x_tx_t^\top
      =\begin{pmatrix}2&0\\0&2\end{pmatrix}
       +\begin{pmatrix}1&2\\2&4\end{pmatrix}
      =\begin{pmatrix}3&2\\2&6\end{pmatrix},\\
    b_2^{\mathrm{nuevo}}
    &=b_2+r_tx_t
      =\begin{pmatrix}0.8\\0.6\end{pmatrix}
       +\begin{pmatrix}1\\2\end{pmatrix}
      =\begin{pmatrix}1.8\\2.6\end{pmatrix}.
\end{align*}

\subsection*{c) Correspondencia con RLHF de un paso}

En la formulación de un paso, cada episodio consta de un prompt,
una respuesta completa y una recompensa escalar. Esto permite
modelarlo como un bandido contextual~\cite{li2010,ouyang2022}.

\begin{itemize}
    \item El contexto $x_t$ corresponde al prompt $x$ que recibe el
    modelo antes de actuar.
    \item El brazo $a_t$ corresponde a la respuesta completa $y$.
    Su generación autorregresiva implementa el muestreo de una sola
    acción en esta formulación.
    \item La política del bandido corresponde a $\pi_\theta(y\mid x)$,
    la distribución condicional del LLM sobre respuestas.
    \item El estimador $\hat r(x_t,a_t)$ corresponde al modelo de
    recompensa $r_\psi(x,y)$ entrenado con comparaciones humanas.
\end{itemize}

Si $p(x)$ es la distribución de prompts, el objetivo es
\[
    \max_\theta\ \mathbb{E}_{x\sim p(x),\,y\sim\pi_\theta(\cdot\mid x)}
    \bigl[R_{\mathrm{RLHF}}(x,y)\bigr],
\]
donde $R_{\mathrm{RLHF}}$ combina el puntaje aprendido con la
regularización respecto a la política SFT. La acción se evalúa
dentro del mismo episodio y la siguiente consulta comienza un
episodio nuevo.
